\documentclass[10pt,a4paper]{amsart}
\usepackage{amsmath,amssymb,amsthm,graphicx}
% Retained minimal preamble for the Stacks Project chapter flat.tex (More on Flatness).
% Verbatim excerpt of the chapter preamble: Xy-pic setup (native preamble lines 14-18) and the native
% theorem declarations with the mathematical macros (native preamble lines 154-217). Copied mechanically;
% nothing added, removed or redefined. The wrapper class and the amsmath/amssymb/amsthm/graphicx base
% packages are supplied by the assembler, so the class/externaldocument part of the native preamble is omitted.
% For commutative diagrams we use Xy-pic
\usepackage[all]{xy}

% We use 2cell for 2-commutative diagrams.
\xyoption{2cell}
% Theorem environments.
%
\theoremstyle{plain}
\newtheorem{theorem}[subsection]{Theorem}
\newtheorem{proposition}[subsection]{Proposition}
\newtheorem{lemma}[subsection]{Lemma}

\theoremstyle{definition}
\newtheorem{definition}[subsection]{Definition}
\newtheorem{example}[subsection]{Example}
\newtheorem{exercise}[subsection]{Exercise}
\newtheorem{situation}[subsection]{Situation}

\theoremstyle{remark}
\newtheorem{remark}[subsection]{Remark}
\newtheorem{remarks}[subsection]{Remarks}

\numberwithin{equation}{subsection}

% Macros
%
\def\lim{\mathop{\mathrm{lim}}\nolimits}
\def\colim{\mathop{\mathrm{colim}}\nolimits}
\def\Spec{\mathop{\mathrm{Spec}}}
\def\Hom{\mathop{\mathrm{Hom}}\nolimits}
\def\Ext{\mathop{\mathrm{Ext}}\nolimits}
\def\SheafHom{\mathop{\mathcal{H}\!\mathit{om}}\nolimits}
\def\SheafExt{\mathop{\mathcal{E}\!\mathit{xt}}\nolimits}
\def\Sch{\mathit{Sch}}
\def\Mor{\mathop{\mathrm{Mor}}\nolimits}
\def\Ob{\mathop{\mathrm{Ob}}\nolimits}
\def\Sh{\mathop{\mathit{Sh}}\nolimits}
\def\NL{\mathop{N\!L}\nolimits}
\def\CH{\mathop{\mathrm{CH}}\nolimits}
\def\proetale{{pro\text{-}\acute{e}tale}}
\def\etale{{\acute{e}tale}}
\def\QCoh{\mathit{QCoh}}
\def\Ker{\mathop{\mathrm{Ker}}}
\def\Im{\mathop{\mathrm{Im}}}
\def\Coker{\mathop{\mathrm{Coker}}}
\def\Coim{\mathop{\mathrm{Coim}}}

% Boxtimes
%
\DeclareMathSymbol{\boxtimes}{\mathbin}{AMSa}{"02}

%
% Macros for moduli stacks/spaces
%
\def\QCohstack{\mathcal{QC}\!\mathit{oh}}
\def\Cohstack{\mathcal{C}\!\mathit{oh}}
\def\Spacesstack{\mathcal{S}\!\mathit{paces}}
\def\Quotfunctor{\mathrm{Quot}}
\def\Hilbfunctor{\mathrm{Hilb}}
\def\Curvesstack{\mathcal{C}\!\mathit{urves}}
\def\Polarizedstack{\mathcal{P}\!\mathit{olarized}}
\def\Complexesstack{\mathcal{C}\!\mathit{omplexes}}
% \Pic is the operator that assigns to X its picard group, usage \Pic(X)
% \Picardstack_{X/B} denotes the Picard stack of X over B
% \Picardfunctor_{X/B} denotes the Picard functor of X over B
\def\Pic{\mathop{\mathrm{Pic}}\nolimits}
\def\Picardstack{\mathcal{P}\!\mathit{ic}}
\def\Picardfunctor{\mathrm{Pic}}
\def\Deformationcategory{\mathcal{D}\!\mathit{ef}}

% --- appended 2026-10-12 (Chief): the project own definitions of the reference/slogan/history
% display environments, copied mechanically from the shared Stacks preamble lines 8-12
% (\usepackage{verbatim} plus the three \newenvironment lines, which the chapter preamble keeps
% verbatim from the project preamble). Nothing added, removed or redefined: these environments
% carry no mathematical content and the project itself typesets them as comments. Without these
% three lines a counted statement whose own text uses them fails with Environment slogan undefined.
\usepackage{verbatim}
\newenvironment{reference}{\comment}{\endcomment}
\newenvironment{slogan}{\comment}{\endcomment}
\newenvironment{history}{\comment}{\endcomment}
\usepackage{cleveref}
\crefname{originaltheorem}{original theorem}{original theorems}
\Crefname{originaltheorem}{Original theorem}{Original theorems}
\makeatletter
\newlabel{algebra-theorem-lazard}{{058G}{source}{}{sourcepointer-1}{}}
\newlabel{algebra-theorem-lazard@cref}{{[originaltheorem][0][]058G}{[1][0][]source}{}{}{}}
\makeatother
\crefname{originaltheorem}{original theorem}{original theorems}
\Crefname{originaltheorem}{Original theorem}{Original theorems}
\makeatletter
\newlabel{algebra-theorem-projective-free-over-local-ring}{{0593}{source}{}{sourcepointer-2}{}}
\newlabel{algebra-theorem-projective-free-over-local-ring@cref}{{[originaltheorem][0][]0593}{[1][0][]source}{}{}{}}
\makeatother
\title{A map of modules over a local ring whose reduction is injective is universally injective}
\author{The Stacks Project Authors}
\date{stacks-project:flat.tex:local-snapshot-20261011; GNU FDL 1.2 or later}
\begin{document}
\maketitle
\noindent\textit{Editorial note.} All delivered mathematical text is the author's own, copied from the pinned Stacks Project chapter file flat.tex; no statement or proof step was rewritten, reordered or removed.\par\medskip
\noindent\textit{Editorial note.} Preamble provenance: the chapter's own retained minimal preamble (sources/stacks-flat/retained-preamble.tex, lines 1-85) is a verbatim excerpt of the native Xy-pic setup and of the native theorem declarations and mathematical macros, followed by the project's own reference/slogan/history display-environment definitions copied mechanically from the shared preamble; the wrapper class and the amsmath/amssymb/amsthm/graphicx base packages are supplied by the assembler, and extra\_packages is left empty because everything needed is already in the retained preamble. No macro was invented or redefined.\par\medskip
\noindent\textit{Editorial note.} Same-chapter pointers are located at their \textbackslash{}label lines inside the pinned chapter file; cross-chapter pointers are located in the project's own tag list staged in the snapshot (sources/stacks-chow/tags.txt), at the verbatim lines carrying their canonical tags.\par\medskip
\noindent\textit{Allowed standard prerequisites (editorial).}\begin{itemize}
\item A projective module over a local ring is free, and Lazard's theorem on flat modules as filtered colimits of finite free modules. Algebra, labels algebra-theorem-projective-free-over-local-ring (proof line 1955) and algebra-theorem-lazard (proof line 1960). The first reduces to the case of a finite free source, and the second writes the target as a directed colimit of finite free modules.
\item Splitting of an injection of finite free modules over a local ring detected modulo the maximal ideal. Algebra, elementary fact used at proof lines 1967-1968. Shows that the reduction of the map being injective implies that the map into each sufficiently late finite free piece is a split injection, which is what makes it stay injective after tensoring with any module.
\end{itemize}\medskip
\ifdefined\hypertarget\hypertarget{sourcepointer-1}{}\fi
\noindent\textit{Source reference.} Original theorem 058G: sources/stacks-chow/tags.txt, source lines 6779--6779. Navigation pointer for \textbackslash{}ref\{algebra-theorem-lazard\} used inside the counted span (line 1960); the argument that consumes it is retained in the span.\par\medskip
\ifdefined\hypertarget\hypertarget{sourcepointer-2}{}\fi
\noindent\textit{Source reference.} Original theorem 0593: sources/stacks-chow/tags.txt, source lines 6801--6801. Navigation pointer for \textbackslash{}ref\{algebra-theorem-projective-free-over-local-ring\} used inside the counted span (line 1955); the argument that consumes it is retained in the span.\par\medskip
\begin{lemma}
\label{lemma-universally-injective-local}
Let $(R, \mathfrak m)$ be a local ring. Let $u : M \to N$ be an $R$-module map.
If $M$ is a projective $R$-module, $N$ is a flat $R$-module, and
$\overline{u} : M/\mathfrak mM \to N/\mathfrak mN$ is injective
then $u$ is universally injective.
\end{lemma}

\begin{proof}
By
Algebra, Theorem \ref{algebra-theorem-projective-free-over-local-ring}
the module $M$ is free. If we show the result holds for every finitely
generated direct summand of $M$, then the lemma follows. Hence we may
assume that $M$ is finite free. Write $N = \colim_i N_i$ as
a directed colimit of finite free modules, see
Algebra, Theorem \ref{algebra-theorem-lazard}.
Note that $u : M \to N$ factors through $N_i$ for some $i$ (as $M$ is finite
free). Denote $u_i : M \to N_i$ the corresponding $R$-module map.
As $\overline{u}$ is injective we see that
$\overline{u_i} : M/\mathfrak mM \to N_i/\mathfrak mN_i$ is
injective and remains injective on composing with the maps
$N_i/\mathfrak mN_i \to N_{i'}/\mathfrak mN_{i'}$ for all $i' \geq i$.
As $M$ and $N_{i'}$ are finite free over the local ring $R$ this implies
that $M \to N_{i'}$ is a split injection for all $i' \geq i$. Hence
for any $R$-module $Q$ we see that $M \otimes_R Q \to N_{i'} \otimes_R Q$
is injective for all $i' \geq i$. As $- \otimes_R Q$ commutes with
colimits we conclude that $M \otimes_R Q \to N_{i'} \otimes_R Q$
is injective as desired.
\end{proof}

\end{document}
